\documentclass[a4paper, top=10mm]{article}
%for writing from the top
\usepackage{fullpage}
%for math
\usepackage{amsmath}
\usepackage{mathrsfs}
\usepackage{amsthm}
%for images
\usepackage{graphicx}
%for color
\usepackage{xcolor}
%for title
\title{\textbf{\huge{Echauffement Geometrique}}}
\author{Enigma n\textsuperscript{o}1}
\date{16\textsuperscript{th} June 2026}

\newtheorem*{hint}{Hint}

\addtolength{\voffset}{-2cm}
\addtolength{\textheight}{5cm}


\begin{document}
	\maketitle
	
	\vspace{1cm}
	
	\Large
	
	Soit un triangle quelconque $ABC$.
	On appelle $M$, $N$, $P$ les pieds de $ABC$ issus respectivement des sommets $A$, $B$, $C$.
	
	\vspace{0.5cm}
	
	Soit $H$ l'orthocentre de $ABC$.
	On appelle $I$, $J$, et $K$ les milieux respectifs des segments $[AC]$, $[BC]$, et $[AH]$.
	
	\vspace{0.5cm}
	
	On cherche la valeur $\angle KMJ \ \times \ \angle JIK$ (multipliez les valeurs en degres).
	
	
	\vspace{3cm}
	
	Consider any triangle $ABC$.
	Let $M$, $N$, and $P$ be the feet of $ABC$ drawn from vertices $A$, $B$, and $C$, respectively.
	
	\vspace{0.5cm}
	
	Let $H$ be the orthocenter of $ABC$.
	Let $I$, $J$, and $K$ be the midpoints of segments $[AC]$, $[BC]$, and $[AH]$, respectively.
	
	\vspace{0.5cm}
	
	We are looking for the value $\angle KMJ \ \times \ \angle JIK$ (multiply both values in degrees).
	
	% Soit C le cercle de diametre [JK], alors I et M apprtiennent a ce cercle, et les angles JIK et JMK valent tous les deux 90 degres, la reponse est donc 8100.
	% Pour M, le fait que M est issu du pied de A dans ABC force un angle de 90 degres entre (MB)=(JM) et (MA)=(MH)=(MK)
	% Pour I, le theoreme des milieu donne (IK) parallele a (CH)=(CP), et (IJ) parallele a (AB)=(AP). P est issus de la hauteur de C, donc (CP) est perpendiculaire a (AP), d'ou (IJ) perpendiculaire a (IJ).
	
\end{document}